% Dodatak A: Šabloni promptova i primeri ocenjivanja
\section{Šabloni promptova i primeri ocenjivanja}
\label{app:promptovi}

Ovaj dodatak sadrži materijale koji podržavaju transparentnost i reproducibilnost sistema EduGrader: osnovne šablone promptova korišćene u eksperimentima i nekoliko reprezentativnih primera ocenjivanja koje je sistem proizveo.
Šabloni promptova opisuju kako se studentski odgovori evaluiraju pod različitim konfiguracijama strogosti.
Prikazani su originalni promptovi korišćeni u eksperimentima.

\medskip
\noindent\textbf{Napomena.} Ocene su prikazane na skali 0--10.

\subsection*{Prompt za blago ocenjivanje}

\begin{tcolorbox}[promptbox]

Ti si veoma blag profesor koji ocenjuje studentske odgovore.

\vspace{2mm}

\textbf{Pravila ocenjivanja:}

\begin{itemize}[leftmargin=1.5em]
\item Cilj je da se prepozna razumevanje koncepta bez obzira na formulaciju odgovora.
\item Nedostajući detalji ili neprecizna terminologija se ne kažnjavaju.
\item Sinonimi, parafraziranje i primeri su prihvatljivi.
\item Kazni samo odgovore koji pokazuju potpuno pogrešno razumevanje ili su van teme.
\item U slučaju nedoumice dodeli više poena.
\item Prazan odgovor dobija 0 poena.
\end{itemize}

\end{tcolorbox}

\subsection*{Podrazumevani prompt (neutralna konfiguracija)}

\begin{tcolorbox}[promptbox]

Ti si profesor na univerzitetu koji ocenjuje studentske odgovore na ispitna pitanja.

\vspace{2mm}

\textbf{Pravila ocenjivanja:}

\begin{itemize}[leftmargin=1.5em]
\item Uporedi studentski odgovor sa tačnim odgovorom i proceni stepen poklapanja.
\item Odgovor ne mora biti identičan tačnom odgovoru ako student pokazuje razumevanje koncepta.
\item Sinonimi, parafraziranje i drugačiji redosled informacija su prihvatljivi ako je suština tačna.
\item Odgovori napisani na različitim jezicima ili pismima (ćirilica/latinica, srpski/engleski) ne treba da budu kažnjeni ako je sadržaj tačan.
\item Ako je odgovor delimično tačan, dodeli proporcionalan broj poena.
\item Ako odgovor sadrži netačne tvrdnje koje nisu deo tačnog odgovora, oduzmi poene proporcionalno.
\end{itemize}

\end{tcolorbox}

\subsection*{Prompt za strogo ocenjivanje}

\begin{tcolorbox}[promptbox]

Ti si veoma strog profesor koji rigorozno ocenjuje studentske odgovore.

\vspace{2mm}

\textbf{Pravila ocenjivanja:}

\begin{itemize}[leftmargin=1.5em]
\item Strogo uporedi studentski odgovor sa tačnim odgovorom.
\item Svaka stavka iz tačnog odgovora koja nedostaje znači oduzimanje poena.
\item Netačne ili neprecizne tvrdnje se kažnjavaju.
\item Površni ili previše opšti odgovori bez konkretnih detalja dobijaju manje poena.
\item Nepotpun odgovor ne može dobiti više od polovine poena.
\item Sinonimi i parafraziranje su prihvatljivi samo ako precizno pokrivaju suštinu tačnog odgovora.
\end{itemize}

\end{tcolorbox}

\subsection*{Reprezentativni primeri ocenjivanja}

Naredni primeri ilustruju tipične izlaze sistema EduGrader.

\subsection*{Primer A: Detekcija stepena broja dva}

\begin{tcolorbox}[promptbox]

\textbf{Pitanje}

Napisati C funkciju koja vraća 1 ako je broj stepen broja dva, a 0 u suprotnom.

\medskip
\textbf{Referentni odgovor}

\begin{verbatim}
int isPowerOfTwo(unsigned x) {
    return x != 0 && (x & (x - 1)) == 0;
}
\end{verbatim}

\medskip
\textbf{Studentski odgovor}

\begin{verbatim}
int isPowerOfTwo(unsigned x) {
    return (x & (x - 1)) == 0;
}
\end{verbatim}

\medskip
\textbf{Ocene}

Nastavnik: 7/10 \hfill Sistem: 8/10

\medskip
\textbf{Obrazloženje sistema}

Student pravilno primenjuje standardnu bitovsku tehniku za detekciju stepena broja dva koristeći izraz \texttt{(x \& (x - 1)) == 0}.
Međutim, rešenje ne obrađuje granični slučaj kada je \texttt{x == 0}.
Pošto izraz vraća tačno za vrednost 0, funkcija daje pogrešan rezultat u tom slučaju.
Rešenje je gotovo tačno, ali nije potpuno.

\end{tcolorbox}

\subsection*{Primer B: Stek i hip memorija}

\begin{tcolorbox}[promptbox]

\textbf{Pitanje}

Koja je razlika između stek memorije i hip memorije u C-u?
Ukratko objasniti.

\medskip
\textbf{Studentski odgovor}

Stek je memorija koju program koristi za pozive funkcija i lokalne promenljive.
Hip je memorija koja čuva globalne promenljive i velike podatke.
Stek je brži, hip je sporiji.

\medskip
\textbf{Ocene}

Nastavnik: 6/10 \hfill Sistem: 6/10

\medskip
\textbf{Obrazloženje sistema}

Student tačno navodi da se stek koristi za pozive funkcija i lokalne promenljive i da je generalno brži.
Međutim, tvrdnja da hip čuva globalne promenljive je netačna.
Globalne promenljive se nalaze u statičkoj memoriji, a ne u hipu.
Takođe, odgovor ne objašnjava dinamičku alokaciju niti razlike u životnom veku memorije.

\end{tcolorbox}
